% SOURCE_AUTHORITY: sga5_fr_workpass.tex, lines 11129--11160 (LF-normalized unit).
% SOURCE_SHA256: 791F4EFFC5E02832D5D77ED03518C8156D6F07E4C8238B03545DB93D883FBB28
\emph{Demonstração.} Como todos os $k$-esquemas considerados são lisos, é imediato, por razões de dimensão, que os morfismos $j_1$ e $\beta$ (resp. $j_1$ e $v$) são transversais. De ([CH] 4 Cor. 1 da prop. 2) resultam então as igualdades
\[
\begin{cases}
j_1^*\beta_*(1)=s_*(1),\\
j_1^*v_*(1)=k_*(1).
\end{cases}
\]
As equações (9.3.2) e (9.3.3) vêm então de (9.1.2) e da normalização das classes de Chern ([CH] 4-18), respectivamente. Demostremos (9.3.4). Como o ciclo $X'$ de $Z'$ é a transformação estrita de $X\times0$ por $f_1$, temos a priori uma igualdade da forma
\[
f_1^*u_*(1)=(j_1)_*g_1^*(\mu)+v_*(1),
\tag{9.3.5}
\]
com $\mu\in\Gamma(Y,\mathbb N)$, e trata-se de ver que $\mu=1$. Isto pode ser feito adaptando a demonstração de (SGA 7 X 2.6.3). Daremos aqui uma demonstração direta. Aplicando $j_1^*$ aos dois membros de (9.3.5), obtém-se (cf. $(D_2)$):
\[
g_1^*i^*u^*u_*(1)=j_1^*(j_1)_*g_1^*(\mu)+j_1^*v_*(1).
\tag{9.3.6}
\]
Um caso muito particular de (9.1.3), ou uma verificação direta, mostra que $u^*u_*(1)=0$. Usando (9.3.3), resulta de (9.3.6) que
\[
0=j_1^*(j_1)_*g_1^*(\mu)+\bar\xi.
\tag{9.3.7}
\]
Restringindo-nos, se necessário, a uma vizinhança de uma componente conexa de $Y$, pode-se supor que $Y$ é conexo, de modo que (9.3.7) toma a forma
\[
0=\mu\,j_1^*(j_1)_*(1)+\bar\xi,\qquad \mu\in\mathbb N.
\tag*{(9.3.7) bis}
\]
Como $j_1^*(j_1)_*(1)$ é homogêneo de grau $1$, é da forma
\[
j_1^*(j_1)_*(1)=g_1^*(\rho)+\nu\bar\xi,
\]
com $\rho\in C^1(Y)$ e $\nu\in\mathbb Z$. Substituindo esta expressão em (9.3.7) bis, obtém-se imediatamente $\rho=0$ e $\mu\nu=-1$, com $\mu\in\mathbb N$, de onde $\mu=1$ e $\nu=-1$, terminando a demonstração de (9.3.4).
